%% =====================================================================
%%  Capítulo 2 -- Modelado I: Laplace, función de transferencia y primer orden
%%  Migrado al formato nuevo (noviembre de 2026) desde Handout2_CtrlParte1.tex
%%  y las presentaciones Modelado1Prof, Modelado2 y Modelado3.
%%  Cambios de fondo respecto a la versión anterior (ver MIGRACION.md):
%%    equilibrio T^* (antes \bar T); escalón \escalon(t) (antes u(t), que
%%    chocaba con la entrada u); polos complejos -1/3 ± j sqrt(2)/3 (faltaba
%%    la j); error ante rampa e(t) = K1 tau (1 - e^{-t/tau}) (signo); se
%%    eliminó la "inversión del tiempo" (no aplica a la transformada
%%    unilateral); el ejemplo de dsolve usa y'(0)=0 en texto y código.
%% =====================================================================

\capitulo[cap:po]{Modelado I: Laplace, función de transferencia y primer orden}
  {La ilusión es el mejor alimento espiritual, el motor del alma. No engaña, entusiasma.}
  {Jorge Barraza}

\begin{objetivos}[partes/02-modelado-primer-orden.html]
  \item Distinguir señal, sistema estático y sistema dinámico, y describir un sistema con un diagrama de bloques de flujo de información.
  \item Escribir el modelo de un proceso a partir de un balance de conservación, con sus suposiciones, y expresarlo en variables desviación.
  \item Calcular transformadas y antitransformadas de Laplace con la tabla, las propiedades y las fracciones parciales, y usarlas para resolver ecuaciones diferenciales lineales.
  \item Definir función de transferencia, polos, ceros, grado relativo y ecuación característica, y usar (bien) el teorema del valor final.
  \item Reconocer un sistema de primer orden, leer $K_1$ y $\tau$ de su respuesta al escalón y predecir su error ante una rampa, verificando todo en \matlab/\simulink.
\end{objetivos}

%% ---------------------------------------------------------------------
\section{Señales y sistemas}
\label{sec:po-senales}
%% ---------------------------------------------------------------------

En el capítulo anterior vimos qué significa ``análisis de sistemas de control'' y cómo nació el área. Ahora empezamos el trabajo de fondo: \emph{construir modelos} de sistemas dinámicos. Todo lo que hagamos en el curso (estabilidad, error, PID, Bode) se apoya en dos palabras que conviene precisar desde ya: señal y sistema.

\begin{definicion}[label=def:po-senal]{Señal}
Una \clave{señal} es una función del tiempo que describe cómo cambia una variable. Se escribe con el tiempo entre paréntesis: $T(t)$, $V(t)$, $y(t)$. Su valor en un instante particular se obtiene evaluándola: $T(2)$ es la temperatura en $t=2$\,s.
\end{definicion}

La \cref{fig:po-senal} muestra la temperatura $T(t)=20-3e^{-0.15t}$ de un recinto que arranca en $17$\,°C y se acerca a $20$\,°C. Esa curva es una señal; la pregunta que nos interesa es \emph{qué la produce}.

\begin{figure}[ht]
\centering
\includegraphics[width=0.8\linewidth]{cap02/f_senal_temperatura}
\caption{Una señal de temperatura. Es la salida de un sistema dinámico de primer orden: en la \cref{sec:po-primerorden} sabremos leer, de esta misma curva, su ganancia y su constante de tiempo.}
\label{fig:po-senal}
\end{figure}

Considere un bombillo dentro de un contenedor aislado (\cref{fig:po-bombillo}). Dos señales nos importan: el voltaje $V(t)$ aplicado al bombillo y la temperatura $T(t)$ dentro del contenedor. Intuimos que a más voltaje, más temperatura: $T(t)$ \emph{depende} de $V(t)$. Esa relación de causa y efecto es lo que llamamos sistema.

\begin{figure}[ht]
\centering
\begin{minipage}{0.36\linewidth}\centering
\includegraphics[width=\linewidth]{cap02/bombillo}
\end{minipage}\hfill
\begin{minipage}{0.6\linewidth}\centering
\includegraphics[width=\linewidth]{cap02/senal_sistema}
\end{minipage}
\caption{Izquierda: el bombillo en un contenedor aislado. Derecha: (a)--(b) el diagrama de bloques de flujo de información: la flecha que entra es la causa, la que sale es el efecto; (c) un sistema dinámico, descrito por una ecuación diferencial; (d) un sistema estático.}
\label{fig:po-bombillo}
\end{figure}

\begin{definicion}[label=def:po-sistema]{Sistema, sistema dinámico}
Un \clave{sistema} es una transformación que toma una señal de entrada $u(t)$ y devuelve una señal de salida $y(t)$. Es \clave{dinámico} cuando la relación entre $u$ e $y$ es una \emph{ecuación diferencial ordinaria} (ODE): involucra a $u$, a $y$ y al menos una de sus derivadas. A esa ecuación la llamamos las \clave{dinámicas} del sistema. Si no hay derivadas, el sistema es \clave{estático}: la salida en $t$ solo depende de la entrada en ese mismo $t$.
\end{definicion}

El bombillo es dinámico (\cref{fig:po-bombillo}c): su temperatura no salta cuando se cambia el voltaje, sino que ``recuerda'' su historia a través de $\dif T/\dif t$. La conversión de unidades del ejemplo siguiente, en cambio, es estática.

\begin{ejemplo}[label=ej:po-fc]{De Fahrenheit a Celsius: un sistema estático}
Si $u(t)$ es una temperatura en °F, la temperatura en °C es $y(t)=\tfrac59\,(u(t)-32)$ (\cref{fig:po-bombillo}d). ¿Cuánto es $90$\,°F?
\solucion $y=\tfrac59(90-32)=\tfrac{290}{9}\approx32.22$\,°C. No hay derivadas: si la entrada cambia de golpe, la salida cambia de golpe. En \matlab se evalúa sobre un vector de temperaturas, numérica o simbólicamente:
\codigorepo{cap02_modelado_primer_orden/C2_Ej1_Senales.m}{fc}{ejemplo \ref{ej:po-fc}}
\end{ejemplo}

\begin{pausa}
La señal de la \cref{fig:po-senal} cumple $\dif T/\dif t+0.15\,T=V$ con $V$ constante. ¿Cuánto vale $V$? ¿Y $T(0)$?
\respuesta{En estado estacionario $\dif T/\dif t=0$, luego $0.15\cdot20=V$, es decir $V=3$. Y $T(0)=20-3=17$\,°C. Se puede comprobar derivando: $\dot T=0.45e^{-0.15t}$ y $\dot T+0.15T=0.45e^{-0.15t}+3-0.45e^{-0.15t}=3$.}
\end{pausa}

%% ---------------------------------------------------------------------
\section{¿Para qué y cómo se modela?}
\label{sec:po-modelar}
%% ---------------------------------------------------------------------

\subsection{El modelo dentro del proceso de diseño}

Un sistema de control se diseña en cuatro grandes pasos: (i) establecer los objetivos, las variables a controlar y las especificaciones con las que se medirá el desempeño; (ii) definir el sistema y \emph{modelarlo} (la configuración, los modelos del proceso, de los actuadores y de los sensores, y su simulación); (iii) escoger el controlador y los parámetros que se van a ajustar; y (iv) optimizar esos parámetros y analizar el desempeño. Este capítulo y el siguiente son el paso (ii). La \cref{fig:po-intercambiador} muestra un caso típico~\citep{SmCo06}: un intercambiador de calor cuya temperatura de salida se controla manipulando el vapor.

\begin{figure}[ht]
\centering
\includegraphics[width=0.5\linewidth]{cap02/intercambiador_pid}\\[-1mm]{\footnotesize (a)}\\[2mm]
\includegraphics[width=\linewidth]{cap02/lazo_proceso}\\[1mm]{\footnotesize (b)}
\caption{(a) Diagrama de instrumentación de un intercambiador de calor (adaptado de~\citep{SmCo06}). (b) El mismo lazo como diagrama de bloques: cada equipo es un bloque y cada flecha, una señal con sus unidades.}
\label{fig:po-intercambiador}
\end{figure}

Antes de escribir una sola ecuación conviene hacerse tres preguntas~\citep{Mar00}:
\begin{itemize}
  \item \textbf{¿Qué variables influyen y cuáles son influenciadas?} Las \clave{variables manipuladas} son las que podemos ajustar (el vapor); las \clave{variables controladas} son las que se ven afectadas (la temperatura de salida). Las demás entradas son perturbaciones.
  \item \textbf{¿En qué rango se pueden mover?} Una válvula no abre más del $100\%$; un calentador no enfría.
  \item \textbf{¿Qué tan efectivamente puede la retroalimentación mantener el proceso donde queremos?} Depende de tres rasgos de la respuesta del proceso: el \emph{signo y la magnitud} (que no cambien con el punto de operación), la \emph{velocidad} (si manipulamos muy agresivamente, el lazo puede oscilar o inestabilizarse) y la \emph{forma} (dos procesos con la misma velocidad pueden ser muy distintos de controlar).
\end{itemize}
Un proceso ``fácil'' tiene ganancia de signo y magnitud estables, responde rápido y su comportamiento no cambia mucho con el tiempo.

\subsection{Modelos orientados al control}

\begin{receta}[Cómo se construye un modelo~\citep{Mar00}]
\begin{enumerate}
  \item \textbf{Definir objetivos:} qué decisiones se van a tomar con el modelo, qué valores numéricos y con qué precisión.
  \item \textbf{Preparar la información:} bosquejo del proceso, variables de interés, \emph{suposiciones} y datos.
  \item \textbf{Formular el modelo:} balances de conservación (masa, energía, momento), ecuaciones constitutivas, y verificar los grados de libertad (tantas ecuaciones como incógnitas).
  \item \textbf{Resolver:} analítica o numéricamente.
  \item \textbf{Analizar los resultados:} ¿tienen sentido? ¿qué tan sensibles son a los parámetros? ¿qué pasa ``si\ldots''?
  \item \textbf{Validar:} comparar con experimentos o con modelos más completos.
\end{enumerate}
\end{receta}

Buscamos \clave{modelos orientados al control} (\emph{control-oriented models}): los más simples que capturen lo esencial de la estática y la dinámica, porque se usan dentro de optimizaciones, en lazos de tiempo real y para diseñar controladores.\nota{``Essentially, all models are wrong, but some are useful'', escribió el estadístico George E.\,P.\ Box~\citep[p.~424]{BoDr87}.} Los modelos de \emph{principios físicos} tienen ventajas sobre los basados solo en datos: explican la dinámica interna, se pueden extrapolar a otras condiciones y se pueden escribir antes de que el equipo exista. Los modelos basados en datos brillan cuando no entendemos bien el proceso, pero su gran debilidad es la extrapolación: nada garantiza que valgan fuera de los datos con que se entrenaron.

\subsection{Sistemas lineales e invariantes en el tiempo}

En casi todo el curso trabajaremos con sistemas \clave{lineales e invariantes en el tiempo} (LTI). Si $y=h(u)$ es la salida para la entrada $u$:
\begin{itemize}
  \item \textbf{Linealidad} (homogeneidad y superposición): $h(a\,u_1+b\,u_2)=a\,h(u_1)+b\,h(u_2)$. Doble entrada, doble salida; la respuesta a una suma es la suma de las respuestas.
  \item \textbf{Invarianza en el tiempo:} si la entrada se retrasa $T$ segundos, la salida se retrasa lo mismo: $u(t-T)\mapsto y(t-T)$. El sistema responde igual hoy que mañana.
\end{itemize}
Ningún sistema real es exactamente LTI.\nota{Richard Feynman lo resumía así: ``Linear systems are important because we can solve them''.} Pero muchos se comportan como tales \emph{alrededor de un punto de operación}, y en esa región podemos usar todo el arsenal de herramientas lineales. En el \cref{cap:so} veremos cómo construir esa aproximación (linealización).

%% ---------------------------------------------------------------------
\subsection{Ejemplo motivador: un tanque agitado}
\label{sec:po-termico-modelo}

Tomemos el tanque de la \cref{fig:po-tanque}~\citep{SmCo06}. Un líquido entra con caudal $f$ y temperatura $T_i(t)$, se mezcla y sale con el mismo caudal y temperatura $T(t)$. Suposiciones: (1) caudales de entrada y salida constantes; (2) densidades constantes; (3) capacidades caloríficas constantes; (4) líquido bien mezclado (la temperatura de salida es la del tanque); (5) tanque bien aislado (sin pérdidas al ambiente); (6) se desprecia la energía que aporta el agitador.

\begin{figure}[ht]
\centering
\includegraphics[width=0.55\linewidth]{cap02/tanque_termico}
\caption{Tanque agitado: ¿cómo responde la temperatura de salida $T(t)$ a un cambio en la de entrada $T_i(t)$? Adaptado de~\citep{SmCo06}.}
\label{fig:po-tanque}
\end{figure}

\textbf{Objetivo:} obtener un modelo matemático y una función de transferencia que describan cómo responde $T(t)$ a un cambio en $T_i(t)$.

El balance de energía (lo que entra menos lo que sale es lo que se acumula) es
\[
  f\rho_i h_i(t)-f\rho\,h(t)=\deriv{}{t}\bigl[V\rho\,u(t)\bigr],
\]
con $f$ el caudal (m$^3$/s), $\rho_i,\rho$ las densidades (kg/m$^3$), $V$ el volumen (m$^3$), $h_i,h$ las entalpías (J/kg) y $u$ la energía interna del líquido (J/kg).\nota{Aquí $u$ es energía interna, no la entrada del sistema. Es la única vez que la usamos así.} Con una temperatura de referencia, $h=C_pT$ y $u=C_vT$; con las suposiciones (2) y (3) queda
\begin{equation}
  f\rho C_p\,T_i(t)-f\rho C_p\,T(t)=V\rho C_v\,\deriv{T(t)}{t}.
  \label{eq:po-balance}
\end{equation}
$T_i(t)$ la imponemos nosotros, así que hay una ecuación y una incógnita: el modelo está completo.

\paragraph{Equilibrio y variables desviación} En estado estacionario la derivada es cero y $f\rho C_p\,\eq{T_i}-f\rho C_p\,\eq{T}=0$: el punto de equilibrio es $\eq{T}=\eq{T_i}$. Definimos las \clave{variables desviación}
\begin{equation}
  \dev{T}(t)=T(t)-\eq{T},\qquad \dev{T}_i(t)=T_i(t)-\eq{T_i}.
  \label{eq:po-desviacion}
\end{equation}
Restando la ecuación de equilibrio de \eqref{eq:po-balance}, y como $\dif\dev{T}/\dif t=\dif T/\dif t$,
\begin{equation}
  \tau\,\deriv{\dev{T}(t)}{t}+\dev{T}(t)=\dev{T}_i(t),\qquad
  \tau=\frac{V C_v}{f C_p}\quad\left[\frac{\text{m}^3\cdot\text{J/(kg\,°C)}}{(\text{m}^3/\text{s})\cdot\text{J/(kg\,°C)}}=\text{s}\right].
  \label{eq:po-termico-ode}
\end{equation}

\begin{ideaclave}
Las variables desviación son un \emph{cambio de coordenadas}: ponen el origen en el punto de equilibrio. Así el sistema ``arranca en cero'', con condiciones iniciales nulas, que es justamente lo que exige la función de transferencia (\cref{sec:po-ft}). Para volver a la variable física basta sumar el equilibrio: $T(t)=\eq{T}+\dev{T}(t)$.
\end{ideaclave}

Podríamos resolver \eqref{eq:po-termico-ode} con cualquier herramienta numérica si conociéramos la entrada. Pero hay una herramienta que además nos da una \emph{representación} del sistema que no depende de la entrada: la transformada de Laplace. La introducimos ahora y volvemos al tanque en la \cref{sec:po-termico}.

%% ---------------------------------------------------------------------
\section{La transformada de Laplace}
\label{sec:po-laplace}
%% ---------------------------------------------------------------------

\begin{historia}[De los operadores de Heaviside a la transformada de Laplace]
Oliver Heaviside (1850--1925) resolvía circuitos y líneas de transmisión con un \emph{cálculo operacional}: reemplazaba $\dif/\dif t$ por un símbolo $p$ y manipulaba las ecuaciones diferenciales como si fueran algebraicas~\citep{Mah17}. Funcionaba, pero sin demostraciones, y los matemáticos de su tiempo lo despreciaron. Thomas Bromwich, en Cambridge, mostró que muchas de esas manipulaciones se podían justificar tratando $p$ como un número complejo e integrando en el plano complejo~\citep{Bro17}; John R.\ Carson, ingeniero de AT\&T, dio en los años veinte una fórmula integral que producía la función operacional a partir de la función del tiempo~\citep{Car26}; y el matemático alemán Gustav Doetsch reunió todo en una teoría rigurosa con el nombre que usamos hoy~\citep{Doe37}, en honor de Pierre-Simon Laplace (1749--1827). El escalón unitario se sigue llamando \emph{función de Heaviside}, y la expansión en fracciones parciales que usaremos para regresar al tiempo también es suya. Heaviside, además, dio a las ecuaciones de Maxwell la forma vectorial que hoy se enseña e hizo posibles las líneas telefónicas largas sin distorsión.
\end{historia}

\begin{definicion}[label=def:po-laplace]{Transformada unilateral de Laplace}
Para una señal $x(t)$ definida para $t\geq0$,
\begin{equation}
  \Lap{x(t)}=X(s)=\int_0^\infty x(t)\,e^{-st}\,\dif t,\qquad s=\sigma+j\omega\in\mathbb{C}.
  \label{eq:po-laplace}
\end{equation}
\end{definicion}

El factor $e^{-st}=e^{-\sigma t}e^{-j\omega t}$ tiene dos partes: una exponencial $e^{-\sigma t}$ que ayuda a que la integral converja y una oscilación $e^{-j\omega t}$. Por eso $s$ es complejo: la transformada descompone la señal en exponenciales y senoides, que son exactamente las funciones que aparecen en las soluciones de las ODE lineales. Con $\sigma=0$ se recupera la transformada de Fourier, que volverá en el capítulo de respuesta frecuencial.

\paragraph{El plano-$s$} Un número complejo $p_1=\sigma_1+j\omega_1$ ($j^2=-1$) tiene parte real $\Re p_1=\sigma_1$, parte imaginaria $\Im p_1=\omega_1$, magnitud $|p_1|=\sqrt{\sigma_1^2+\omega_1^2}$ y fase $\angle p_1=\phi_1$ (el ángulo medido desde el semieje real positivo; ojo con el cuadrante, $\tan^{-1}(\omega_1/\sigma_1)$ sola no lo distingue). Su conjugado es $\bar p_1=\sigma_1-j\omega_1$ y $|p_1|^2=p_1\bar p_1$. Todos se dibujan en el \clave{plano complejo-$s$} (\cref{fig:po-planos}).

\begin{figure}[ht]
\centering
\includegraphics[width=0.62\linewidth]{cap02/plano_s}
\caption{Un número complejo y su conjugado en el plano-$s$. La fase se mide desde el semieje real positivo: aquí $\phi_1$ está en el segundo cuadrante.}
\label{fig:po-planos}
\end{figure}

\paragraph{Existencia} Basta que $x(t)$ sea \emph{continua a trozos} (en cada intervalo finito tiene un número finito de saltos, con límites a izquierda y derecha) y de \emph{orden exponencial} (existe $a$ tal que $e^{-at}|x(t)|$ es acotada para $t$ grande). La segunda condición impone $\Re s>\sigma_a$ para algún $\sigma_a$: es la \emph{región de convergencia}. En el curso todas las señales cumplen estas condiciones y no insistiremos en la región de convergencia.

\begin{cuidado}[Todas las señales empiezan en $t=0$]
La transformada unilateral solo ``ve'' $t\geq0$. Por eso, cuando escribimos $x(t)$ queremos decir $x(t)\,\escalon(t)$, donde $\escalon(t)$ es el escalón unitario (1 para $t\geq0$ y 0 antes). En las notas anteriores el escalón se escribía $u(t)$; aquí lo escribimos $\escalon(t)$ porque $u(t)$ es la \emph{entrada} de un sistema.
\end{cuidado}

\begin{ejemplo}[label=ej:po-lap-escalon]{Escalón y exponencial, desde la definición}
Calcule $\Lap{\escalon(t)}$ y $\Lap{e^{-at}}$ con $a>0$.
\solucion
\[
  \Lap{\escalon(t)}=\int_0^\infty e^{-st}\,\dif t=\Bigl[-\tfrac1s e^{-st}\Bigr]_0^\infty=\frac1s\quad(\Re s>0),
\]
porque $e^{-st}\to0$ cuando $\sigma>0$. De igual forma,
\[
  \Lap{e^{-at}}=\int_0^\infty e^{-(s+a)t}\,\dif t=\frac{1}{s+a}\quad(\Re s>-a).
\]
\end{ejemplo}

\Cref{tab:po-transformadas} reúne los pares que usaremos todo el semestre. \matlab los calcula con el \emph{Symbolic Math Toolbox}:
\codigorepo{cap02_modelado_primer_orden/C2_Ej2_Laplace.m}{transformadas}{tabla \ref{tab:po-transformadas}}

\begin{table}[ht]
\centering
\caption{Pares de transformadas de Laplace ($t\geq0$, $a,\omega>0$).}
\label{tab:po-transformadas}
\renewcommand{\arraystretch}{2.3}
\begin{tabular}{@{}>{$}l<{$}>{$}l<{$}@{\qquad}>{$}l<{$}>{$}l<{$}@{}}
\toprule
x(t) & X(s) & x(t) & X(s)\\
\midrule
\impulso(t) & 1 & t\,e^{-at} & \dfrac{1}{(s+a)^2}\\
\escalon(t) & \dfrac1s & t^n e^{-at} & \dfrac{n!}{(s+a)^{n+1}}\\
t & \dfrac1{s^2} & \sin\omega t & \dfrac{\omega}{s^2+\omega^2}\\
t^n & \dfrac{n!}{s^{n+1}} & \cos\omega t & \dfrac{s}{s^2+\omega^2}\\
e^{-at} & \dfrac{1}{s+a} & e^{-at}\sin\omega t & \dfrac{\omega}{(s+a)^2+\omega^2}\\
 & & e^{-at}\cos\omega t & \dfrac{s+a}{(s+a)^2+\omega^2}\\
\bottomrule
\end{tabular}
\end{table}

%% ---------------------------------------------------------------------
\subsection{Propiedades}
\label{sec:po-propiedades}

Con unas pocas propiedades la tabla alcanza para casi todo. Sea $x(t)\leftrightarrow X(s)$.

\begin{enumerate}
  \item \textbf{Linealidad:} $a\,x_1(t)+b\,x_2(t)\ \leftrightarrow\ a\,X_1(s)+b\,X_2(s)$. (La integral es lineal.)
  \item \textbf{Corrimiento en el tiempo} ($t_0>0$): $x(t-t_0)\,\escalon(t-t_0)\ \leftrightarrow\ e^{-st_0}X(s)$.
  \item \textbf{Corrimiento en $s$:} $e^{s_0t}x(t)\ \leftrightarrow\ X(s-s_0)$. Por ejemplo, $e^{-2t}\cos3t\leftrightarrow\dfrac{s+2}{(s+2)^2+9}$, y como $\cos\omega_0t=\tfrac12(e^{j\omega_0t}+e^{-j\omega_0t})$, $\cos(\omega_0t)\,x(t)\leftrightarrow\tfrac12X(s-j\omega_0)+\tfrac12X(s+j\omega_0)$.
  \item \textbf{Escalamiento} ($a>0$): $x(at)\ \leftrightarrow\ \dfrac1a X\!\left(\dfrac sa\right)$. (Con el cambio de variable $\tau=at$.)
  \item \textbf{Diferenciación:} $\dot x(t)\ \leftrightarrow\ sX(s)-x(0)$ y, en general,
  \begin{equation}
    \frac{\dif^n x}{\dif t^n}\ \leftrightarrow\ s^nX(s)-s^{n-1}x(0)-s^{n-2}\dot x(0)-\cdots-x^{(n-1)}(0).
    \label{eq:po-derivada}
  \end{equation}
  \item \textbf{Integración:} $\displaystyle\int_0^t x(\lambda)\,\dif\lambda\ \leftrightarrow\ \frac{X(s)}{s}$.
  \item \textbf{Diferenciación en $s$:} $t\,x(t)\ \leftrightarrow\ -\dfrac{\dif X(s)}{\dif s}$. (De aquí sale $t\,e^{-at}\leftrightarrow1/(s+a)^2$.)
  \item \textbf{Convolución:} $\displaystyle (x_1*x_2)(t)=\int_0^t x_1(\lambda)\,x_2(t-\lambda)\,\dif\lambda\ \leftrightarrow\ X_1(s)\,X_2(s)$.
\end{enumerate}

La propiedad de diferenciación es la que convierte una ecuación diferencial en una algebraica. Su demostración es una integración por partes con $u=x(t)$ y $\dif v=e^{-st}\dif t$:
\[
  \Lap{\dot x}=\int_0^\infty \dot x\,e^{-st}\,\dif t
  =\Bigl[x\,e^{-st}\Bigr]_0^\infty+s\int_0^\infty x\,e^{-st}\,\dif t
  =-x(0)+sX(s),
\]
porque $x\,e^{-st}\to0$ (orden exponencial). Aplicándola dos veces se obtiene la de la segunda derivada, y así sucesivamente. Observe que \alerta{las condiciones iniciales aparecen solas}: no hay que agregarlas al final como con los métodos clásicos.

\codigorepo{cap02_modelado_primer_orden/C2_Ej2_Laplace.m}{propiedades}{propiedades de la transformada}

\begin{profundiza}[¿Y la ``inversión del tiempo''?]
La versión anterior de estas notas listaba la propiedad $x(-t)\leftrightarrow X(-s)$. Es válida para la transformada \emph{bilateral} (integral de $-\infty$ a $\infty$), pero no para la unilateral que usamos: si $x(t)$ solo existe para $t\geq0$, entonces $x(-t)$ solo existe para $t\leq0$ y la integral \eqref{eq:po-laplace} la ve como cero. Por la misma razón el escalamiento solo se enuncia con $a>0$.
\end{profundiza}

\begin{pausa}
Use las propiedades y la tabla para obtener $\Lap{3-2e^{-4t}+5t\,e^{-4t}}$.
\respuesta{Por linealidad, $\dfrac3s-\dfrac{2}{s+4}+\dfrac{5}{(s+4)^2}$.}
\end{pausa}

%% ---------------------------------------------------------------------
\section{Resolver ecuaciones diferenciales con Laplace}
\label{sec:po-odes}
%% ---------------------------------------------------------------------

El método tiene tres pasos: \textbf{(1)} transformar la ODE con \eqref{eq:po-derivada}, incluyendo las condiciones iniciales; \textbf{(2)} despejar $X(s)$, que queda como un cociente de polinomios; \textbf{(3)} regresar al tiempo expandiendo $X(s)$ en fracciones parciales y usando la tabla.

\begin{ejemplo}[label=ej:po-motor]{El calibrador óptico}
Un calibrador óptico que verificaba la alineación de las piezas metálicas de varios sedanes de cuatro puertas de Chrysler tenía dinámicas no forzadas (la ODE igualada a cero)
\[
  \ddot x+5\dot x+6x=0,\qquad x(0)=1,\ \dot x(0)=0.
\]
Halle $X(s)$ y $x(t)$.
\solucion Transformando término a término,
\[
  \bigl[s^2X-s\,x(0)-\dot x(0)\bigr]+5\bigl[sX-x(0)\bigr]+6X=0
  \;\Longrightarrow\;
  X(s)(s^2+5s+6)=s+5,
\]
\[
  X(s)=\frac{s+5}{(s+2)(s+3)}=\frac{3}{s+2}-\frac{2}{s+3}
  \quad\Longrightarrow\quad x(t)=3e^{-2t}-2e^{-3t}.
\]
La \cref{fig:po-motor} muestra la respuesta y sus dos componentes.
\end{ejemplo}

\begin{figure}[ht]
\centering
\includegraphics[width=\linewidth]{cap02/f_chrysler}
\caption{Ejemplo~\ref{ej:po-motor}. Cada polo de $X(s)$ aporta un \emph{modo} $e^{p_it}$; el residuo es su peso. El polo en $-3$ muere más rápido que el de $-2$.}
\label{fig:po-motor}
\end{figure}

\codigorepo{cap02_modelado_primer_orden/C2_Ej3_Motor.m}{simbolico}{ejemplo \ref{ej:po-motor} con Laplace simbólico}

\begin{definicion}[label=def:po-polos]{Polos y ceros}
Si $X(s)=N(s)/D(s)$ es un cociente de polinomios, los \clave{ceros} de $X$ son las raíces de $N(s)=0$ y los \clave{polos} son las raíces de $D(s)=0$. En el plano-$s$ se dibujan los polos con $\times$ y los ceros con $\circ$.
\end{definicion}

La expansión en fracciones parciales, que introdujo Heaviside, depende del tipo de polos. Supondremos $X(s)$ estrictamente propia (grado de $N$ menor que el de $D$); si no lo es, se divide primero.

\subsection{Polos simples}
Si $D(s)=(s-p_1)(s-p_2)\cdots(s-p_n)$ con raíces distintas,
\begin{equation}
  X(s)=\sum_{i=1}^n\frac{k_i}{s-p_i},\qquad
  k_i=\Bigl[(s-p_i)\,X(s)\Bigr]_{s=p_i},\qquad
  x(t)=\sum_{i=1}^n k_i\,e^{p_it}.
  \label{eq:po-residuos}
\end{equation}
En el ejemplo~\ref{ej:po-motor}: $k_1=\bigl[(s+5)/(s+3)\bigr]_{s=-2}=3$ y $k_2=\bigl[(s+5)/(s+2)\bigr]_{s=-3}=-2$. A los $k_i$ se les llama \clave{residuos}, y \matlab los calcula con \cmd{residue}:
\codigorepo{cap02_modelado_primer_orden/C2_Ej3_Motor.m}{residue}{residuos, polos y diagrama de polos y ceros}

\subsection{Polos múltiples}
Si $p$ es un polo de multiplicidad $r$, aporta $r$ términos:
\[
  \frac{k_{1}}{s-p}+\frac{k_{2}}{(s-p)^2}+\cdots+\frac{k_{r}}{(s-p)^r},\qquad
  k_{r-m}=\frac{1}{m!}\,\frac{\dif^m}{\dif s^m}\Bigl[(s-p)^rX(s)\Bigr]_{s=p},\ m=0,\dots,r-1,
\]
y en el tiempo, por la tabla, $k_m\dfrac{t^{m-1}}{(m-1)!}e^{pt}$: aparecen términos $t\,e^{pt}$, $t^2e^{pt}$\ldots

\begin{ejemplo}[label=ej:po-multiples]{Dos polos dobles}
Expanda $X(s)=\dfrac{1}{s^2(s+1)^2}$ y obtenga $x(t)$.
\solucion Con $s=0$ (doble): $k_2=\bigl[1/(s+1)^2\bigr]_0=1$ y $k_1=\bigl[\tfrac{\dif}{\dif s}(s+1)^{-2}\bigr]_0=-2$. Con $s=-1$ (doble): $k_2=\bigl[1/s^2\bigr]_{-1}=1$ y $k_1=\bigl[-2s^{-3}\bigr]_{-1}=2$. Luego
\[
  X(s)=-\frac2s+\frac1{s^2}+\frac2{s+1}+\frac1{(s+1)^2},\qquad
  x(t)=-2+t+2e^{-t}+t\,e^{-t}.
\]
\end{ejemplo}
\codigorepo{cap02_modelado_primer_orden/C2_Ej4_FraccionesParciales.m}{multiples}{ejemplo \ref{ej:po-multiples}}

\subsection{Polos complejos conjugados}
Un par $-\alpha\pm j\beta$ tiene residuos complejos conjugados y produce, juntos, un término real $e^{-\alpha t}(A\cos\beta t+B\sin\beta t)$. En la práctica es más cómodo \emph{completar el cuadrado} en el denominador y usar las dos últimas filas de la \cref{tab:po-transformadas}.

\begin{ejemplo}[label=ej:po-complejos]{Una oscilación amortiguada}
Resuelva $3\ddot x+2\dot x+x=0$ con $x(0)=-1$, $\dot x(0)=2$.
\solucion Transformando, $3(s^2X+s-2)+2(sX+1)+X=0$, de donde
\[
  X(s)=\frac{4-3s}{3s^2+2s+1},\qquad\text{polos: } s=-\tfrac13\pm j\,\tfrac{\sqrt2}{3}.
\]
Completando el cuadrado, $3s^2+2s+1=3\bigl[(s+\tfrac13)^2+\tfrac29\bigr]$, y escribiendo el numerador en términos de $s+\tfrac13$:
\[
  X(s)=\frac{-(s+\frac13)+\frac53}{(s+\frac13)^2+(\frac{\sqrt2}{3})^2}
  \;\Longrightarrow\;
  x(t)=e^{-t/3}\left[-\cos\frac{\sqrt2\,t}{3}+\frac{5\sqrt2}{2}\sin\frac{\sqrt2\,t}{3}\right].
\]
Se comprueba que $x(0)=-1$ y $\dot x(0)=\tfrac13+\tfrac53=2$. La parte real de los polos ($-\tfrac13$) da la envolvente; la imaginaria ($\tfrac{\sqrt2}3$), la frecuencia de oscilación (\cref{fig:po-complejos}).
\end{ejemplo}

\begin{figure}[ht]
\centering
\includegraphics[width=0.85\linewidth]{cap02/f_polos_complejos}
\caption{Ejemplo~\ref{ej:po-complejos}: la parte real de los polos fija la envolvente $e^{-t/3}$; la imaginaria, la frecuencia $\sqrt2/3$\,rad/s.}
\label{fig:po-complejos}
\end{figure}
\codigorepo{cap02_modelado_primer_orden/C2_Ej4_FraccionesParciales.m}{complejos}{ejemplo \ref{ej:po-complejos}}

\subsection{De los polos a la forma de la respuesta}
Los tres casos se resumen en un solo mapa (\cref{fig:po-modos}): cada polo $p$ aporta un modo $e^{pt}$. La parte real dice si el modo crece o muere, y qué tan rápido; la parte imaginaria dice si oscila, y a qué frecuencia. Un polo en el semiplano derecho basta para que la respuesta crezca sin límite.

\begin{figure}[ht]
\centering
\includegraphics[width=\linewidth]{cap02/f_polos_modos}
\caption{Cada ubicación de polo en el plano-$s$ produce un modo distinto. Izquierda del eje imaginario: el modo se extingue (más rápido cuanto más a la izquierda). Sobre el eje: oscila sin amortiguarse. A la derecha: crece.}
\label{fig:po-modos}
\end{figure}

\begin{ideaclave}
Los polos determinan \emph{qué} modos aparecen en la respuesta; los residuos (y por lo tanto los ceros y las condiciones iniciales) determinan \emph{cuánto} pesa cada uno. Volveremos sobre esto en el \cref{cap:so} con los polos dominantes.
\end{ideaclave}

\subsection{Sin transformar: \texttt{dsolve} y \texttt{ode45}}
\matlab puede resolver el problema de valor inicial directamente: \cmd{dsolve} lo hace simbólicamente (internamente, con las mismas ideas) y \cmd{ode45} numéricamente, sin necesidad de que la ecuación sea lineal.
\codigorepo{cap02_modelado_primer_orden/C2_Ej5_dsolve.m}{dsolve}{$\dot y=-2y+t$ y $\ddot y+\dot y+y=\sin t$ con $y(0)=1$, $\dot y(0)=0$}

\begin{pausa}
En el ejemplo~\ref{ej:po-motor}, ¿qué condición inicial $\dot x(0)$ haría desaparecer el modo $e^{-3t}$?
\respuesta{Con $\dot x(0)=v$, $X(s)=\dfrac{s+5+v}{(s+2)(s+3)}$. El residuo en $-3$ es $\dfrac{-3+5+v}{-1}=-(2+v)$, que se anula con $v=-2$. Entonces $X=1/(s+2)$ y $x(t)=e^{-2t}$: la condición inicial ``cae'' justo sobre el modo lento.}
\end{pausa}

%% ---------------------------------------------------------------------
\section{Función de transferencia}
\label{sec:po-ft}
%% ---------------------------------------------------------------------

Las herramientas anteriores nos llevan al concepto central de todo el curso.

\begin{definicion}[label=def:po-ft]{Función de transferencia}
La \clave{función de transferencia} de un sistema LTI es el cociente entre la transformada de Laplace de la salida y la de la entrada, \emph{con condiciones iniciales nulas}:
\begin{equation}
  T(s)=\left.\frac{\Lap{\text{salida}}}{\Lap{\text{entrada}}}\right|_{\text{CI}=0}=\frac{Y(s)}{U(s)}.
\end{equation}
\end{definicion}

Si el sistema se describe con la ecuación entrada-salida
\begin{equation}
  \frac{\dif^n y}{\dif t^n}+a_1\frac{\dif^{n-1}y}{\dif t^{n-1}}+\cdots+a_n y
  =b_0\frac{\dif^m u}{\dif t^m}+\cdots+b_{m-1}\frac{\dif u}{\dif t}+b_m u,
  \label{eq:po-inout}
\end{equation}
la linealidad y la diferenciación con condiciones iniciales nulas dan
\begin{equation}
  T(s)=\frac{Y(s)}{U(s)}=\frac{b_0s^m+b_1s^{m-1}+\cdots+b_m}{s^n+a_1s^{n-1}+\cdots+a_n}.
  \label{eq:po-ft}
\end{equation}

Algunos nombres que usaremos todo el tiempo:
\begin{itemize}
  \item $T(s)$ es \clave{propia} si $n\geq m$, y \clave{estrictamente propia} si $n>m$. Los sistemas físicos son propios: una función impropia derivaría la entrada y amplificaría sin límite el ruido de alta frecuencia.
  \item El \clave{grado relativo} es $n-m$. El \clave{orden} del sistema es $n$.
  \item El polinomio del denominador igualado a cero, $s^n+a_1s^{n-1}+\cdots+a_n=0$, es la \clave{ecuación característica}. Sus raíces son los polos de $T(s)$: los modos naturales del sistema.
\end{itemize}

La función de transferencia permite describir el sistema con un bloque (\cref{fig:po-bloqueft}): la señal de salida es el producto $Y(s)=T(s)\,U(s)$. En el tiempo, por la propiedad de convolución, eso es $y(t)=\int_0^t g(\lambda)\,u(t-\lambda)\,\dif\lambda$, donde $g(t)=\iLap{T(s)}$ es la \clave{respuesta al impulso}: el producto en $s$ es mucho más cómodo que la integral en $t$.

\begin{figure}[ht]
\centering
\includegraphics[width=0.62\linewidth]{cap02/bloque_ft}
\caption{Un bloque de flujo de información: la función de transferencia es un mapeo lineal entre la transformada de la entrada y la de la salida. Las reglas para combinar bloques están en el \cref{cap:so}.}
\label{fig:po-bloqueft}
\end{figure}

\begin{cuidado}[La función de transferencia no ve las condiciones iniciales]
$T(s)$ describe la respuesta \emph{forzada} desde el reposo. Si el sistema no arranca en reposo, la respuesta total es la forzada más la libre (la que producen las condiciones iniciales, como en el ejemplo~\ref{ej:po-motor}). Por eso trabajamos en variables desviación alrededor de un equilibrio: allí las condiciones iniciales son cero.
\end{cuidado}

\begin{pausa}
Para $\dddot y+2\ddot y+3\dot y+4y=\dddot u+5\dot u+6u$, escriba $T(s)$. ¿Es propia? ¿Cuál es su grado relativo?
\respuesta{$T(s)=\dfrac{s^3+5s+6}{s^3+2s^2+3s+4}$. Es propia pero no estrictamente propia ($n=m=3$): grado relativo $0$. Una parte de la entrada pasa directamente a la salida.}
\end{pausa}

%% ---------------------------------------------------------------------
\section{El teorema del valor final}
\label{sec:po-tvf}
%% ---------------------------------------------------------------------

Muchas veces no necesitamos toda la respuesta, solo a dónde llega: el \clave{valor en estado estacionario}. Para eso existe una de las propiedades más útiles de la transformada.

\begin{definicion}[label=def:po-tvf]{Teorema del valor final}
Si \alerta{todos los polos de $sY(s)$ tienen parte real negativa}, entonces
\begin{equation}
  \lim_{t\to\infty}y(t)=\lim_{s\to0}sY(s).
  \label{eq:po-tvf}
\end{equation}
\end{definicion}

La condición es esencial: el teorema supone que el límite en el tiempo existe. Si no existe, la fórmula de la derecha da un número igual, y ese número no significa nada.

\begin{ejemplo}[label=ej:po-tvf]{Dos casos en los que el teorema miente}
Determine $\lim_{t\to\infty}h(t)$ para (a) $H(s)=\dfrac1{s-2}$ y (b) $H(s)=\dfrac{6(s^2-2)}{s^3+s^2-4s-4}$.
\solucion (a) Aplicando \eqref{eq:po-tvf} a ciegas, $\lim_{s\to0}\frac{s}{s-2}=0$. Pero $h(t)=e^{2t}$ crece sin cota: el polo de $sH(s)$ en $+2$ invalida el teorema.

(b) De nuevo $\lim_{s\to0}sH(s)=0$. Pero $s^3+s^2-4s-4=(s+1)(s+2)(s-2)$, y
\[
  H(s)=\frac{2}{s+1}+\frac{3}{s+2}+\frac{1}{s-2}\;\Longrightarrow\;h(t)=2e^{-t}+3e^{-2t}+e^{2t}\to\infty.
\]
Moraleja: \emph{primero} los polos, \emph{después} el límite.
\end{ejemplo}

\codigorepo{cap02_modelado_primer_orden/C2_Ej6_ValorFinal.m}{tvf}{ejemplo \ref{ej:po-tvf}: revisar los polos antes del límite}

\begin{profundiza}[El teorema del valor inicial]
Con una condición parecida (que el límite exista), $\displaystyle\lim_{t\to0^+}y(t)=\lim_{s\to\infty}sY(s)$. Sirve, por ejemplo, para saber si una respuesta al escalón arranca con salto (grado relativo $0$) o arranca en cero con pendiente (grado relativo $1$): aplíquelo a $sY(s)$ para obtener $\dot y(0^+)$.
\end{profundiza}

%% ---------------------------------------------------------------------
\section{De vuelta al tanque: la primera función de transferencia}
\label{sec:po-termico}
%% ---------------------------------------------------------------------

\begin{ejemplo}[label=ej:po-termico]{El tanque agitado ante un escalón}
El modelo \eqref{eq:po-termico-ode} es $\tau\,\dif\dev{T}/\dif t+\dev{T}=\dev{T}_i$. (a) Obtenga la función de transferencia. (b) Si la temperatura de entrada, que estaba en el equilibrio $\eq{T_i}=\eq{T}$, sube bruscamente $M$ grados, ¿cómo responde la salida? (c) ¿A qué valor llega?
\solucion (a) Con condiciones iniciales nulas (variables desviación), $\tau s\,\Dev{T}(s)+\Dev{T}(s)=\Dev{T}_i(s)$, luego
\begin{equation}
  \frac{\Dev{T}(s)}{\Dev{T}_i(s)}=\frac{1}{\tau s+1}.
  \label{eq:po-termico-ft}
\end{equation}
(b) La entrada es $\dev{T}_i(t)=M\,\escalon(t)$, es decir, $\Dev{T}_i(s)=M/s$. Por polos simples,
\[
  \Dev{T}(s)=\frac{M}{s(\tau s+1)}=\frac{M}{s}-\frac{M}{s+1/\tau}
  \;\Longrightarrow\;
  \dev{T}(t)=M\left(1-e^{-t/\tau}\right),\qquad
  T(t)=\eq{T}+M\left(1-e^{-t/\tau}\right).
\]
(c) Todos los polos de $s\Dev{T}(s)=M/(\tau s+1)$ están en el semiplano izquierdo, así que el teorema aplica: $\lim_{s\to0}s\Dev{T}(s)=M$. Es lo mismo que se obtiene de $\lim_{t\to\infty}\dev{T}(t)$. La salida sube exactamente lo que subió la entrada: la ganancia del tanque es $1$, como pide la conservación de la energía sin pérdidas.
\end{ejemplo}

\begin{figure}[ht]
\centering
\includegraphics[width=\linewidth]{cap02/f_termico}
\caption{Ejemplo~\ref{ej:po-termico} con $\eq{T}=\eq{T_i}=25$\,°C, $M=2$\,°C y $\tau=5$\,s. Izquierda: las variables desviación, que arrancan en cero. Derecha: las temperaturas físicas, que se obtienen sumando el equilibrio.}
\label{fig:po-termico}
\end{figure}

\codigorepo{cap02_modelado_primer_orden/C2_Ej7_Termico.m}{termico}{ejemplo \ref{ej:po-termico}}

\begin{cuidado}[\texttt{step} es un escalón \emph{unitario}]
\cmd{step(G)} aplica $\escalon(t)$. Para un escalón de tamaño $M$ se multiplica: \cmd{M*step(G,t)} (por linealidad). Y \cmd{step} devuelve la \emph{desviación}: si quiere la temperatura física, sume el equilibrio.
\end{cuidado}

\begin{modelo}[Termico.slx]
El script \cmd{crearModelo\_Termico.m} construye por código el diagrama escalón (o rampa) $\to$ $K_1/(\tau s+1)$ $\to$ $+\eq{T}$, con un \emph{Scope} para las variables desviación y otro para las físicas. Cambie \cmd{tau}, \cmd{M} y \cmd{Teq} en el \emph{workspace} y vuelva a simular; con el interruptor manual pase a la entrada rampa de la \cref{sec:po-primerorden}.
\end{modelo}

\begin{pausa}
Si se duplica el caudal $f$ del tanque, ¿la respuesta es más rápida o más lenta? ¿Cambia el valor final?
\respuesta{Más rápida: $\tau=VC_v/(fC_p)$ se reduce a la mitad (el líquido se renueva en la mitad del tiempo). El valor final no cambia: la ganancia sigue siendo $1$.}
\end{pausa}

%% ---------------------------------------------------------------------
\section{El sistema de primer orden}
\label{sec:po-primerorden}
%% ---------------------------------------------------------------------

El tanque es el ejemplo típico de un \clave{sistema de primer orden}:
\begin{equation}
  \frac{Y(s)}{U(s)}=\frac{K_1}{\tau s+1},
  \label{eq:po-primerorden}
\end{equation}
con dos parámetros: la \clave{ganancia} $K_1$ (cuánto cambia la salida en estado estacionario por unidad de cambio en la entrada) y la \clave{constante de tiempo} $\tau$ (qué tan rápido). El único polo está en $s=-1/\tau$. Veamos qué hacen ante las dos entradas de prueba más usadas.

\subsection{Respuesta al escalón unitario}
Con $U(s)=1/s$,
\begin{equation}
  Y(s)=\frac{K_1}{s}-\frac{K_1}{s+1/\tau}\quad\Longrightarrow\quad y(t)=K_1\left(1-e^{-t/\tau}\right).
  \label{eq:po-escalon}
\end{equation}
La respuesta tiene dos partes: la de estado estacionario $y_{ss}=K_1$ (el teorema del valor final da lo mismo) y la transitoria $-K_1e^{-t/\tau}$, que se extingue. La \cref{fig:po-escalon} muestra los números que vale la pena memorizar.

\begin{figure}[ht]
\centering
\includegraphics[width=0.92\linewidth]{cap02/f_primer_orden_escalon}
\caption{Respuesta al escalón de $K_1/(\tau s+1)$ con $K_1=10$ y $\tau=5$. Cada constante de tiempo, la salida recorre el $63.2\%$ de lo que le falta. La tangente en el origen llega al valor final justo en $t=\tau$.}
\label{fig:po-escalon}
\end{figure}

\begin{itemize}
  \item En $t=\tau$ la salida llega al $1-e^{-1}=63.2\%$ del valor final; en $2\tau$, al $86.5\%$; en $3\tau$, al $95.0\%$; en $4\tau$, al $98.2\%$; en $5\tau$, al $99.3\%$.
  \item \clave{Tiempo de establecimiento} (2\%): $\ts\approx4\tau$ (exactamente $\tau\ln50=3.91\,\tau$).
  \item La pendiente inicial es $\dot y(0)=K_1/\tau$: si la respuesta siguiera con esa pendiente, llegaría al valor final en $t=\tau$. Constante de tiempo pequeña, respuesta rápida.
\end{itemize}

\codigorepo{cap02_modelado_primer_orden/C2_Ej8_PrimerOrden.m}{escalon}{escalón: 63.2 \%, $\ts$ y \cmd{stepinfo}}

\begin{receta}[Identificar un primer orden con una prueba al escalón]
\begin{enumerate}
  \item Llevar el proceso a un equilibrio $(\eq{u},\eq{y})$ y aplicar un escalón $\Delta u$.
  \item Esperar el nuevo estado estacionario y medir $\Delta y$: $K_1=\Delta y/\Delta u$ (con sus unidades).
  \item Buscar el instante en que la salida recorrió el $63.2\%$ de $\Delta y$: ese tiempo es $\tau$. (O, con un solo dato $y(t_1)$: $\tau=-t_1/\ln\bigl(1-\dev{y}(t_1)/\Delta y\bigr)$.)
  \item Escribir $\Dev{Y}/\Dev{U}=K_1/(\tau s+1)$ y verificar superponiendo la simulación a los datos.
\end{enumerate}
\end{receta}

\subsection{Respuesta a la rampa}
Con $u(t)=t$, es decir $U(s)=1/s^2$,
\[
  Y(s)=\frac{K_1}{s^2}-\frac{K_1\tau}{s}+\frac{K_1\tau}{s+1/\tau}
  \;\Longrightarrow\;
  y(t)=K_1t-K_1\tau+K_1\tau\,e^{-t/\tau},\quad t\geq0.
\]
La parte $K_1t-K_1\tau$ es la de estado estacionario y $K_1\tau e^{-t/\tau}$ la transitoria. Definiendo el error de seguimiento como $e(t)=K_1u(t)-y(t)$ (lo que la salida ``debería'' valer si no tuviera dinámica),
\begin{equation}
  e(t)=K_1\tau\left(1-e^{-t/\tau}\right)\;\xrightarrow{\ t\to\infty\ }\;K_1\tau.
  \label{eq:po-error-rampa}
\end{equation}
La salida sigue a la rampa con la misma pendiente, pero \emph{atrasada} $\tau$ segundos: el error final es $K_1\tau$ (\cref{fig:po-rampa}). Cuanto menor $K_1\tau$, mejor el seguimiento.

\begin{figure}[ht]
\centering
\includegraphics[width=\linewidth]{cap02/f_primer_orden_rampa}
\caption{Respuesta a la rampa con $K_1=1$: con $\tau=10$ el error final es $10$; con $\tau=1$, es $1$.}
\label{fig:po-rampa}
\end{figure}

\codigorepo{cap02_modelado_primer_orden/C2_Ej8_PrimerOrden.m}{rampa}{rampa: el error tiende a $K_1\tau$}

\begin{cuidado}[El signo del error ante la rampa]
En la versión anterior de estas notas aparecía $e(t)=K_1\tau+K_1\tau e^{-t/\tau}$, que vale $2K_1\tau$ en $t=0$. Lo correcto es $e(t)=K_1\tau(1-e^{-t/\tau})$: en $t=0$ el error es nulo (la rampa arranca en cero) y crece hasta $K_1\tau$.
\end{cuidado}

\begin{pausa}
Un primer orden con $K_1=2$ responde a un escalón de $3$ unidades y a los $12$\,s ya está en el $86.5\%$ del cambio final. ¿Cuánto valen $\tau$, el cambio final y el $\ts$?
\respuesta{$86.5\%$ corresponde a $2\tau$, luego $\tau=6$\,s. Cambio final $K_1\cdot3=6$ unidades; $\ts\approx4\tau=24$\,s.}
\end{pausa}

%% ---------------------------------------------------------------------
\section{Para llevar}
\label{sec:po-cierre}
%% ---------------------------------------------------------------------

\begin{parallevar}
\begin{enumerate}
  \item Un sistema transforma una señal de entrada en una de salida. Es dinámico si la relación es una ODE: la salida ``recuerda''.
  \item Un modelo orientado al control parte de un balance con suposiciones explícitas. En el equilibrio se definen las variables desviación $\dev{x}=x-\eq{x}$: arrancan en cero.
  \item Laplace convierte derivadas en productos por $s$ \emph{e incluye las condiciones iniciales}. Para regresar al tiempo: fracciones parciales (polos simples, múltiples o complejos) y la tabla.
  \item Cada polo aporta un modo $e^{pt}$: parte real, crecimiento o decaimiento; parte imaginaria, oscilación. Los residuos dicen cuánto pesa cada modo.
  \item Función de transferencia: $Y(s)/U(s)$ con condiciones iniciales nulas. Propia, grado relativo, ecuación característica.
  \item Teorema del valor final: \emph{primero} verificar que los polos de $sY(s)$ estén en el semiplano izquierdo.
  \item Primer orden $K_1/(\tau s+1)$: $63.2\%$ en $\tau$, $\ts\approx4\tau$, pendiente inicial $K_1/\tau$, error ante rampa $K_1\tau$.
\end{enumerate}
\end{parallevar}

\begin{profundiza}[Correcciones respecto a la versión anterior del capítulo]
(i) El punto de equilibrio se escribe $\eq{T}$ (antes $\bar T$) y el escalón $\escalon(t)$ (antes $u(t)$, que chocaba con la entrada). (ii) Polos complejos del ejemplo~\ref{ej:po-complejos}: $-\tfrac13\pm j\tfrac{\sqrt2}{3}$ (faltaba la $j$). (iii) La fase de un número complejo es $\angle p$ (el texto decía ``magnitud''). (iv) Se eliminó la ``inversión del tiempo'', que no vale para la transformada unilateral. (v) El ejemplo de \cmd{dsolve} usa $\dot y(0)=0$ en el texto y en el código. (vi) El error ante la rampa es $K_1\tau(1-e^{-t/\tau})$. (vii) Se quitó una mención a ``la posición final de la masa'' que no correspondía a ningún ejemplo. Todos los números del capítulo se verificaron con \texttt{herramientas/cap02\_numeros.py} y \cmd{verificarCap02.m}.
\end{profundiza}

%% ---------------------------------------------------------------------
\section{Problemas}
\label{sec:po-problemas}
%% ---------------------------------------------------------------------

Los problemas marcados con \etiquetaM{} requieren \matlab/\simulink; los marcados con \etiquetaEstrella{} son de profundización.

\subsection*{Transformadas y fracciones parciales}

\begin{problema}[transformadas]\label{prob:po-1}
Determine $X(s)$ para: (a) $x(t)=3-2e^{-4t}+5te^{-4t}$; (b) $x(t)=e^{-2t}\cos3t+4t^2$; (c) $x(t)=t\,e^{-t}\sin2t$; (d) $x(t)=\int_0^t e^{-3\lambda}\,\dif\lambda$; (e) el pulso $x(t)=2$ para $0\leq t<3$ y $0$ en otro caso. En (a)--(d) use la tabla y las propiedades; en (e), la definición \eqref{eq:po-laplace}, y diga qué propiedad le habría permitido evitar la integral. \etiquetaM{} Verifique con \cmd{laplace}.
\end{problema}

\begin{problema}[fracciones parciales]\label{prob:po-2}
Para $F_1(s)=\dfrac{s+3}{(s+1)(s+4)}$, $F_2(s)=\dfrac{10}{s(s+2)^2}$ y $F_3(s)=\dfrac{2s+5}{s^2+4s+13}$: (a) ubique polos y ceros en el plano-$s$; (b) expanda en fracciones parciales y obtenga $f_i(t)$; (c) \etiquetaM{} verifique con \cmd{residue} y \cmd{pzplot}; (d) diga qué términos desaparecen cuando $t\to\infty$ y relaciónelo con la ubicación de los polos.
\end{problema}

\begin{problema}[ODE con condiciones iniciales]\label{prob:po-3}
Considere $2\ddot y+10\dot y+12y=24u$ con $u(t)=\escalon(t)$, $y(0)=1$, $\dot y(0)=0$. (a) Obtenga $Y(s)$ sin despreciar las condiciones iniciales. (b) Halle $y(t)$ y verifique las dos condiciones iniciales. (c) Calcule $\lim_{t\to\infty}y(t)$ de dos maneras. (d) Repita (c) para $\ddot y-\dot y-2y=3u$ con condiciones iniciales nulas; obtenga $y(t)$ y explique por qué el teorema del valor final falla. (e) \etiquetaM{} Resuelva ambos con \cmd{ode45} y grafique.
\end{problema}

\subsection*{Función de transferencia}

\begin{problema}[función de transferencia]\label{prob:po-4}
Para (i) $\dddot y+2\ddot y+3\dot y+4y=\dddot u+5\dot u+6u$; (ii) $\ddot y+7\dot y+10y=4\dot u+8u$; (iii) $5\dot y+2y=3u$: (a) escriba $T(s)$; (b) diga si es propia o estrictamente propia y su grado relativo; (c) halle polos y ceros; (d) ¿cuál es de primer orden de la forma \eqref{eq:po-primerorden}? Identifique $K_1$ y $\tau$. (e) Mire con cuidado el sistema (ii): ¿qué le pasa a $T(s)$?
\end{problema}

\begin{problema}[del bloque a la ODE]\label{prob:po-5}
Escriba la ecuación diferencial que representa el bloque de la \cref{fig:po-repaso}.
\end{problema}

\begin{figure}[ht]
\centering
\includegraphics[width=0.6\linewidth]{cap02/bloque_repaso}
\caption{Problema~\ref{prob:po-5}.}
\label{fig:po-repaso}
\end{figure}

\subsection*{Primer orden}

\begin{problema}[exploración]\label{prob:po-6}
\etiquetaM{} Con \cmd{C2\_Ej8\_PrimerOrden.m} y \cmd{crearModelo\_Termico.m}, varíe $K_1$ y $\tau$ y compruebe las afirmaciones de la \cref{sec:po-primerorden} para el escalón y la rampa. ¿Qué parámetro cambia el valor final y cuál la velocidad?
\end{problema}

\begin{problema}[identificación de un proceso térmico]\label{prob:po-7}
Un proceso opera en estado estacionario con temperatura de salida $\eq{T}=40$\,°C y entrada $\eq{u}=10$ unidades. En $t=0$ la entrada pasa a $15$ unidades; la salida se estabiliza en $52$\,°C y a los $30$\,s marca $47.6$\,°C.
(a) Halle $K_1$ y $\tau$ y escriba la función de transferencia en variables desviación.
(b) Calcule la pendiente inicial y el $\ts$ (2\%).
(c) ¿Qué entrada lleva la salida a $55$\,°C?
(d) Si desde el equilibrio la entrada sube como una rampa de $0.01$ unidades/s, halle $y(t)$ y el error final $\lim e(t)$ con $e=K_1\dev{u}-\dev{y}$.
(e) Si ese error no puede superar $0.5$\,°C, ¿cuál es el máximo $\tau$ admisible? ¿Qué implicaría para el diseño del equipo?
(f) \etiquetaM{} Verifique (a)--(d) en \simulink.
\end{problema}

\begin{problema}[teorema del valor inicial]\label{prob:po-8}
\etiquetaEstrella{} Use el teorema del valor inicial para hallar $y(0^+)$ y $\dot y(0^+)$ ante un escalón unitario en (a) $\dfrac{K_1}{\tau s+1}$ y (b) $\dfrac{s+3}{s+1}$. Relacione el resultado con el grado relativo y explique por qué en (b) la salida ``salta''.
\end{problema}

%% ---------------------------------------------------------------------
\section{Soluciones seleccionadas}
%% ---------------------------------------------------------------------

\paragraph{Problema~\ref{prob:po-1}}
(a) $\dfrac3s-\dfrac2{s+4}+\dfrac5{(s+4)^2}$. (b) $\dfrac{s+2}{(s+2)^2+9}+\dfrac8{s^3}$. (c) $\Lap{e^{-t}\sin2t}=\dfrac{2}{(s+1)^2+4}$ y, por diferenciación en $s$, $\dfrac{4(s+1)}{\bigl[(s+1)^2+4\bigr]^2}$. (d) Por integración, $\dfrac{1}{s(s+3)}$. (e) $\int_0^3 2e^{-st}\dif t=\dfrac{2(1-e^{-3s})}{s}$; es $2\escalon(t)-2\escalon(t-3)$ y basta el corrimiento en el tiempo.

\paragraph{Problema~\ref{prob:po-2}}
$f_1=\tfrac23e^{-t}+\tfrac13e^{-4t}$ (polos $-1,-4$; cero $-3$). $f_2=2.5-2.5e^{-2t}-5t\,e^{-2t}$ (polos $0$ y $-2$ doble): el término constante no desaparece porque hay un polo en el origen. $f_3=e^{-2t}\bigl(2\cos3t+\tfrac13\sin3t\bigr)$ (polos $-2\pm3j$, cero $-2.5$).

\paragraph{Problema~\ref{prob:po-3}}
(a) $Y(s)=\dfrac{s^2+5s+12}{s(s+2)(s+3)}$. (b) $y(t)=2-3e^{-2t}+2e^{-3t}$; $y(0)=1$, $\dot y(0)=6-6=0$. (c) $2$ por los dos caminos. (d) $Y=\dfrac{3}{s(s-2)(s+1)}$, el teorema daría $-1.5$; pero $y(t)=-1.5+0.5e^{2t}+e^{-t}$ diverge por el polo en $+2$.

\paragraph{Problema~\ref{prob:po-4}}
(i) $\dfrac{s^3+5s+6}{s^3+2s^2+3s+4}$, propia, grado relativo $0$. (ii) $\dfrac{4(s+2)}{(s+2)(s+5)}$: el cero cancela el polo en $-2$ y queda $\dfrac{4}{s+5}=\dfrac{0.8}{0.2s+1}$; el modo $e^{-2t}$ no se ve desde la entrada (volveremos sobre esto en espacio de estado). (iii) $\dfrac{3}{5s+2}=\dfrac{1.5}{2.5s+1}$: primer orden con $K_1=1.5$ y $\tau=2.5$.

\paragraph{Problema~\ref{prob:po-5}}
$y^{(4)}+5\dddot y+6\ddot y+7\dot y+8y=\dddot u+2\ddot u+3\dot u+4u$.

\paragraph{Problema~\ref{prob:po-7}}
(a) $K_1=12/5=2.4$\,°C/unidad; $1-e^{-30/\tau}=7.6/12$ da $\tau=29.9\approx30$\,s; $\Dev{T}/\Dev{U}=2.4/(29.9s+1)$. (b) $0.40$\,°C/s y $\ts\approx4\tau\approx120$\,s. (c) $\dev{u}=15/2.4=6.25$, es decir $u=16.25$. (d) $\dev{y}=0.024\bigl(t-\tau+\tau e^{-t/\tau}\bigr)$ y $e_\infty=K_1a\tau=0.72$\,°C. (e) $\tau\leq0.5/(2.4\cdot0.01)=20.8$\,s: habría que hacer el proceso un $30\%$ más rápido (menos volumen o más caudal), o, mejor, cerrar un lazo de control.

\script[verificación automática de los números del capítulo]{cap02_modelado_primer_orden/verificarCap02.m}

\respuestaspausas

